% LaTeX master Datei(en) zusammengestellt von Falk-Jonatan Strube zur Nutzung an der Hochschule für Technik und Wirtschaft Dresden: https://github.com/genericFJS/htw
\documentclass{scrreprt}
\gdef\pathtomaster{../_LaTeX_master}
%% HTW Dresden Corporate Design
%% Copyright (C) Falk-Jonatan Strube <falkjonatanstrube@gmail.com>, 2017
%%% === Anfang Paket htwscr ===
\makeatletter
\ifcsname KOMAoption\endcsname%
\KOMAoption{footheight}{11mm}	% mehr Platz für Wasserzeichen o.ä. in der Fußzeile
\recalctypearea	% s.o.
\RequirePackage{geometry}
\geometry{top=25mm, bottom=25mm, headsep=10mm, inner=25mm, outer=25mm, footskip=15mm, marginparsep=2mm}
\RequirePackage{helvet}	% Corporate Design schreibt Arial/Helvetica vor
\renewcommand{\familydefault}{\sfdefault}	% s.o.
\RequirePackage{graphicx}	% für Grafiken
\RequirePackage{xcolor}		% für Farben
\definecolor{htworange}{RGB}{249,155,28}	% Corporate Design orange CMYK(0/45/100/0)
\definecolor{htwblue}{RGB}{0,116,188}			% Corporate Design blau CMYK(90/50/0/0)
\definecolor{htwgrey}{RGB}{128,128,128}		% Corporate Design gray 60% Schwarz [Annahme: CMYK(0,0,0,60)]
\RequirePackage{ifthen,xifthen}	% für Bedingungen/Optionen
\RequirePackage{etoolbox}				% s.o.

%%% Format des HTW-Headers (wird hier als Texteelement verwenden) nach dem Style der PowerPoint-Vorlagen. Weiterhin inspiriert durch das Paket tudscr (Corporate Design der TU Dresden)
%%% Dateipfad des Logos abhängig von der Verwendung als Paket/Input
\ifdefined\@isHTWCDpackage%
	\gdef\@pathtoimage{HTW}%
\else%
	\ifdefined\pathtomaster%
		\gdef\@pathtoimage{\pathtomaster /HTW}%
	\else 
		\gdef\@pathtoimage{HTW}%
	\fi%
\fi%
\newcommand{\pathtologo}[1]{\gdef\@pathtoimage{#1}}
%%% Berechnung der Größe des Logos
\newlength{\htwLogoHeight}
\setlength{\htwLogoHeight}{12mm}	% Mindesthöhe des Logos laut Corporate Design: 11mm
\newlength{\@htwLogoWidth}
\AtBeginDocument{
\def\printLogo{\includegraphics[height=\htwLogoHeight]{\@pathtoimage}}
\setlength\@htwLogoWidth{\widthof{\printLogo}}
}
%%% Definition des Header
\newcommand{\@htwheader}{{\noindent%
%%% Header ist abhängig von dem Vorhandensein von department, institute und chair
\newcommand{\hasDepInsCh}{\gdef\@hasDepInsCh{}}%
\ifdefined\@department\hasDepInsCh\fi%
\ifdefined\@institute\hasDepInsCh\fi%
\ifdefined\@chair\hasDepInsCh\fi%
\begin{minipage}[b][\htwLogoHeight]{\dimexpr \textwidth - \@htwLogoWidth - 2mm \relax}% 2mm Abstand zwischen Logo und Linie
{\small{\bfseries%
\ifdefined\@faculty\@faculty{}\fi} %
\strut}%
{\color{htworange}\hrule height 1pt}%
\vspace{.1mm}%
\end{minipage}
\hfill%
\begin{minipage}[b][\htwLogoHeight]{\@htwLogoWidth}
\includegraphics[height=\htwLogoHeight]{\@pathtoimage}
\end{minipage}
\ifdefined\@hasDepInsCh%
\\%
\begin{minipage}{\textwidth}
\small%
\ifdefined\@department\@department{}\fi%
\ifdefined\@institute\ifdefined\@department,~\fi \@institute\fi%
\ifdefined\@chair\ifdefined\@department,~\else\ifdefined\@institute,~\fi\fi\@chair\fi\strut
%{\color{htworange}\hrule height 0.5pt}
\end{minipage}
\fi
}}

%%% Definition der Fußzeilen
\usepackage{scrlayer-scrpage}
\@ifclassloaded{scrreprt}{%
\renewcommand*{\partpagestyle}{scrheadings}
\renewcommand*{\chapterpagestyle}{scrheadings}%
}{}
\pagestyle{scrheadings}
\cfoot{\vfill\pagemark}
\chead{}

%%% Einstellungen
\newcommand{\simpledate}{\gdef\@simpledate{}}		% Datum nicht mit "Eingereicht am:" Vorsatz
\newcommand{\nouppercase}{\gdef\@nouppercase{}}	% Titel/Parts/Chapters usw. nicht in Großbuchstaben

%%% Einiges erst zu Beginn des Dokumentes ausführen, um Einstellungen zu berücksichtigen
\AtBeginDocument{
\renewcommand\SS{SS}	% fix big ß
%%% Sektionen usw. ggf. in Großbuchstaben:
\providecommand{\chapterlinesformat}[3]{}	% compatibility
\providecommand{\chapterlineswithprefixformat}[3]{}	% compatibility
\providecommand{\sectionlinesformat}[4]{}	% compatibility
\ifdefined\@nouppercase\relax\else
\renewcommand\sectionlinesformat[4]{\@hangfrom{\hskip#2 #3}{\MakeUppercase{#4}}}
\@ifclassloaded{scrreprt}{%
\renewcommand\chapterlinesformat[3]{\@hangfrom{#2}{\MakeUppercase{#3}}}
\renewcommand\chapterlineswithprefixformat[3]{\MakeUppercase{#2#3}}%
}{}
\fi
%%% Anpassung der Part-Seite
\renewcommand*{\partheadstartvskip}{\@htwheader\vspace*{30mm}}
\renewcommand*{\raggedpart}{\raggedright}
\renewcommand*{\partformat}{\huge\ifdefined\@nouppercase \partname~\thepart\else\MakeUppercase{\partname~\thepart}\fi.}
\ifdefined\@nouppercase\setkomafont{part}{\Huge}\else\setkomafont{part}{\Huge\MakeUppercase}\fi
}

\newcommand{\bilingual}[3]{\ifdefined #1 \let #1\undefined\fi\newcommand{#1}{\iflanguage{english}{#3}{#2}}}	% Übersetzung von einigen Variablen in mehrere Sprachen

\newcommand{\checkThesis}[2]{\ifthenelse{\equal{#1}{diss}\OR\equal{#1}{doctoral}\OR\equal{#1}{phd}}{
\gdef#2{\dissertationname}}{
\ifthenelse{\equal{#1}{diploma}}{
\gdef#2{\diplomathesisname}}{
\ifthenelse{\equal{#1}{master}}{
\gdef#2{\masterthesisname}}{
\ifthenelse{\equal{#1}{bachelor}}{
\gdef#2{\bachelorthesisname}}{
\ifthenelse{\equal{#1}{student}}{
\gdef#2{\studentthesisname}}{
\ifthenelse{\equal{#1}{evidence}}{
\gdef#2{\studentresearchname}}{
\ifthenelse{\equal{#1}{project}}{
\gdef#2{\projectpapername}}{
\ifthenelse{\equal{#1}{seminar}}{
\gdef#2{\seminarpapername}}{
\ifthenelse{\equal{#1}{term}}{
\gdef#2{\termpapername}}{
\ifthenelse{\equal{#1}{research}}{
\gdef#2{\researchname}}{
\ifthenelse{\equal{#1}{log}}{
\gdef#2{\logname}}{
\ifthenelse{\equal{#1}{report}}{
\gdef#2{\reportname}}{
\ifthenelse{\equal{#1}{internship}}{
\gdef#2{\internshipname}}{
\ifthenelse{\equal{#1}{lecture}\OR\equal{#1}{lesson}\OR\equal{#1}{pract}}{
\bilingual{\professorname}{Vorlesung von}{Lecture by}
\simpledate
\gdef\@islecture{\relax}
\ifthenelse{\equal{#1}{lecture}}{\gdef#2{\lecturename}}{
\ifthenelse{\equal{#1}{lesson}}{\gdef#2{\lessonname}}{
\gdef#2{\practname}
}}}{
\gdef#2{#1}}}}}}}}}}}}}}}}

%%% Textfragmente der Titelseite in deutsch und englisch
\bilingual{\dateofbirthname}{Geboren am:}{Born on:}
\bilingual{\placeofbirthname}{in}{in}
\bilingual{\coursename}{Studiengang:}{Course:}
\bilingual{\disciplinename}{Studienrichtung:}{Discipline:}
\bilingual{\matriculationnumbername}{Matrikelnummer:}{Matriculation number:}
\bilingual{\matriculationyearname}{Immatrikulationsjahr}{Matriculation year:}
\bilingual{\graduationname}{zur Erlangung des akademischen Grades}{to achieve the academic degree}
\bilingual{\refereename}{Gutachter}{Referee}
\bilingual{\advisorname}{Fachreferent}{Advisor}
\bilingual{\supervisorname}{Betreuer}{Supervisor}
\bilingual{\professorname}{Betreuender Hochschullehrer}{Supervising professor}
\bilingual{\datename}{Eingereicht am:}{Submitted on:}
%% Thesis Typen (inspiriert von dem Paket tudscr):
\bilingual{\dissertationname}{Dissertation}{Dissertation} % diss, doctoral, phd 
\bilingual{\diplomathesisname}{Diplomarbeit}{Diploma Thesis} % diploma 
\bilingual{\masterthesisname}{Master-Arbeit}{Master Thesis} % master 
\bilingual{\bachelorthesisname}{Bachelor-Arbeit}{Bachelor Thesis} % bachelor 
\bilingual{\studentthesisname}{Studienarbeit}{Student Thesis} % student 
\bilingual{\studentresearchname}{Großer Beleg}{Student Research Project} % evidence 
\bilingual{\projectpapername}{Projektarbeit}{Project Paper} % project
\bilingual{\seminarpapername}{Seminararbeit}{Seminar Paper} % seminar
\bilingual{\termpapername}{Hausarbeit}{Term Paper} % term
\bilingual{\researchname}{Forschungsbericht}{Research Report} % research 
\bilingual{\logname}{Protokoll}{Log} % log 
\bilingual{\reportname}{Bericht}{Report} % report
\bilingual{\internshipname}{Praktikumsbericht}{Internship Report} % internship 
\bilingual{\lecturename}{Vorlesungsmitschrift}{Lecture Notes} % lecture 
\bilingual{\lectureauthorname}{Mitschrift von}{Notes by} % lecture  
\bilingual{\lessonname}{\"Ubungsmitschrift}{Lesson Notes} % lecture  
\bilingual{\practname}{Praktikumsmitschrift}{Pracital Course Notes} % lecture 
%%% Variablen für die Titelseite (inspiriert von dem Paket tudscr):
\renewcommand{\subject}[1]{\checkThesis{#1}{\@subject}}
\newcommand{\faculty}[1]{\gdef\@faculty{#1}}
\newcommand{\department}[1]{\gdef\@department{#1}}
\newcommand{\institute}[1]{\gdef\@institute{#1}}
\newcommand{\chair}[1]{\gdef\@chair{#1}}
\newcommand{\authormore}[1]{\gdef\@authormore{#1}}
\newcommand{\dateofbirth}[1]{\gdef\@dateofbirth{#1}}
\newcommand{\placeofbirth}[1]{\gdef\@placeofbirth{#1}}
\newcommand{\course}[1]{\gdef\@course{#1}}
\newcommand{\discipline}[1]{\gdef\@discipline{#1}}
\newcommand{\matriculationnumber}[1]{\gdef\@matriculationnumber{#1}}
\newcommand{\matriculationyear}[1]{\gdef\@matriculationyear{#1}}
\newcommand{\thesis}[1]{\checkThesis{#1}{\@thesis}}
\newcommand{\graduation}[2][]{\gdef\@graduation{#2}\gdef\@graduationshort{#1}}
\newcommand{\supervisor}[1]{\gdef\@supervisor{#1}}
\newcommand{\referee}[1]{\gdef\@referee{#1}}
\newcommand{\advisor}[1]{\gdef\@advisor{#1}}
\newcommand{\professor}[1]{\gdef\@professor{#1}}
%%% Titelseite (inspiriert von dem Paket tudscr):
\renewcommand*{\maketitle}{%
\begin{titlepage}%
\begin{flushleft}
\renewcommand{\thefootnote}{$\star$}%
%%% Kopfzeile HTW-Style:
\@htwheader\\[13mm]
%%% Die Zeilen für titlehead, subject, title, subtitle und author sind immer vorhanden (spacing)
{\noindent\@titlehead{}\strut}\\[3mm]%
{\Large\@subject{}\strut}\\[8mm]%
{\Huge\bfseries%
\ifdefined\@nouppercase\@title{}\else\MakeUppercase{\@title{}}\fi%
\strut}\\
{\Large\bfseries%
\ifdefined\@nouppercase\@subtitle{}\else\MakeUppercase{\@subtitle{}}\fi%
\strut}\\[11mm]
\ifdefined\@islecture\strut\else{\Large\@author{}\strut}\fi % Autor wird bei Vorlesungsmitschrift unten angezeigt
%%% Weitere Felder werden nur angezeigt, wenn sie definiert wurden
\ifdefined\@authormore\ifdefined\@islecture\else\\\fi{\@authormore{}\strut}\fi
\ifdefined\@dateofbirth\\\dateofbirthname{} \@dateofbirth{} \ifdefined\@placeofbirth\placeofbirthname{} \@placeofbirth{}\strut\fi\fi
\ifdefined\@course\\\coursename{} \@course{}\strut\fi
\ifdefined\@discipline\\\disciplinename{} \@discipline{}\strut\fi
\ifdefined\@matriculationnumber\\\matriculationnumbername{} \@matriculationnumber{}\strut\fi
\ifdefined\@matriculationyear\\\matriculationyearname{} \@matriculationyear{}\strut\fi
\ifdefined\@thesis\\[10mm]{\LARGE\bfseries
\ifdefined\@nouppercase\@thesis{}\else\MakeUppercase{\@thesis{}}\fi
\strut}
\ifdefined\@graduation\\[2mm]{\graduationname{}\strut}
\\[5mm]{\LARGE\bfseries
\ifdefined\@nouppercase\relax\else\MakeUppercase{\fi%
\@graduation{} \ifdefined\@graduationshort(\@graduationshort{})\fi
\ifdefined\@nouppercase\relax\else}\fi\strut}\fi\fi
\vfill\noindent
\ifdefined\@islecture\\[1mm]{\lectureauthorname{}\strut}\\[.7mm]{\Large\@author{}\strut}\fi
\ifdefined\@referee{\refereename{}\strut}\\[.7mm]{\Large\@referee{}\strut}\fi
\ifdefined\@advisor\\[1mm]{\advisorname{}\strut}\\[.7mm]{\Large\@advisor{}\strut}\fi
\ifdefined\@supervisor\\[1mm]{\supervisorname{}\strut}\\[.7mm]{\Large\@supervisor{}\strut}\fi
\ifdefined\@professor\\[1mm]{\professorname{}\strut}\\[.7mm]{\Large\@professor{}\strut}\fi
\ifdefempty{\@date}{}{\ifdefined\@simpledate\\[6mm]\relax\else\\[3mm]\datename{}~\fi\@date{}\strut}
\ifdefempty{\@publishers}{}{\\[4.5mm]{\Large\@publishers{}\strut}}%
\renewcommand\footnoterule{}
\@thanks\global\let\@thanks\@empty
\end{flushleft}
\end{titlepage}
\setcounter{footnote}{0}
}
\else%
\PackageError{htwcd}{htwcd needs a KOMA class!}%
\fi%
\makeatother
%%% === Ende Paket htwscr ===
%-- Nummerierung bis Subsubsection für Report
\setcounter{secnumdepth}{3}	% zählt auch subsubsection
\setcounter{tocdepth}{3}	% Inhaltsverzeichnis bis in subsubsection

% LuaLaTeX Unterscheidung
\usepackage{ifluatex}

% Input inkl. Umlaute, Silbentrennung
\ifluatex
	\usepackage{fontspec}
\else
	\usepackage[utf8]{inputenc}
	\usepackage[T1]{fontenc}
\fi
\usepackage[english,ngerman]{babel}
\usepackage{csquotes}	% Anführungszeichen
\usepackage{marvosym}
\usepackage{eurosym}

% Style-Aufhübschung
\usepackage{soul, color}	% Kapitälchen, Unterstrichen, Durchgestrichen usw. im Text
%\usepackage{titleref}
\usepackage[breakwords, fit]{truncate}	% Abschneiden von Sätzen
\renewcommand{\TruncateMarker}{\,…}

% Mathe usw.
\usepackage{amssymb}
\usepackage{amsthm}
\usepackage[fleqn,intlimits]{amsmath}	% fleqn: align-Umgebung rechtsbündig; intlimits: Integralgrenzen immer ober-/unterhalb
%\usepackage{mathtools} % u.a. schönere underbraces
\usepackage{xcolor}
\usepackage{esint}	% Schönere Integrale, \oiint vorhanden
\everymath=\expandafter{\the\everymath\displaystyle}	% Mathe Inhalte werden weniger verkleinert
\usepackage{wasysym}	% mehr Symbole, bspw \lightning
%\renewcommand{\int}{\int\limits}
%\usepackage{xfrac}	% mehr fracs: sfrac{}{}
\let\oldemptyset\emptyset	% schöneres emptyset
\let\emptyset\varnothing
%\RequirePackage{mathabx}	% mehr Symbole
\mathchardef\mhyphen="2D	% Hyphen in Math
\usepackage{mathtools}		% zum Beispiel {pmatrix}[r] für Formatierung der Matrix

% tikz usw.
\usepackage{tikz}
\usepackage{pgfplots}
\pgfplotsset{compat=1.11}	% Umgeht Fehlermeldung
\usetikzlibrary{graphs}
%\usetikzlibrary{through}	% ???
\usetikzlibrary{arrows}
\usetikzlibrary{arrows.meta}	% Pfeile verändern / vergrößern: \draw[-{>[scale=1.5]}] (-3,5) -> (-3,3);
\usetikzlibrary{automata,positioning} % Zeilenumbruch im Node node[align=center] {Text\\nächste Zeile} automata für Graphen
\usetikzlibrary{matrix}
\usetikzlibrary{patterns}	% Schraffierte Füllung
\usetikzlibrary{shapes.geometric}	% Polygon usw.
\tikzstyle{reverseclip}=[insert path={	% Inverser Clip \clip
	(current page.north east) --
	(current page.south east) --
	(current page.south west) --
	(current page.north west) --
	(current page.north east)}
% Nutzen: 
%\begin{tikzpicture}[remember picture]
%\begin{scope}
%\begin{pgfinterruptboundingbox}
%\draw [clip] DIE FLÄCHE, IN DER OBJEKT NICHT ERSCHEINEN SOLL [reverseclip];
%\end{pgfinterruptboundingbox}
%\draw DAS OBJEKT;
%\end{scope}
%\end{tikzpicture}
]	% Achtung: dafür muss doppelt kompliert werden!
\usepackage{graphpap}	% Grid für Graphen
\tikzset{every state/.style={inner sep=2pt, minimum size=2em}}
\usetikzlibrary{mindmap, backgrounds}
%\usepackage{tikz-uml}	% braucht Dateien: http://perso.ensta-paristech.fr/~kielbasi/tikzuml/

% Tabular
\usepackage{longtable}	% Große Tabellen über mehrere Seiten
\usepackage{multirow}	% Multirow/-column: \multirow{2[Anzahl der Zeilen]}{*[Format]}{Test[Inhalt]} oder \multicolumn{7[Anzahl der Reihen]}{|c|[Format]}{Test2[Inhalt]}
\renewcommand{\arraystretch}{1.3} % Tabellenlinien nicht zu dicht
\usepackage{colortbl}
\arrayrulecolor{gray}	% heller Tabellenlinien
\usepackage{array}	% für folgende 3 Zeilen (für Spalten fester breite mit entsprechender Ausrichtung):
\newcolumntype{L}[1]{>{\raggedright\let\newline\\\arraybackslash\hspace{0pt}}m{\dimexpr#1\columnwidth-2\tabcolsep-1.5\arrayrulewidth}}
\newcolumntype{C}[1]{>{\centering\let\newline\\\arraybackslash\hspace{0pt}}m{\dimexpr#1\columnwidth-2\tabcolsep-1.5\arrayrulewidth}}
\newcolumntype{R}[1]{>{\raggedleft\let\newline\\\arraybackslash\hspace{0pt}}m{\dimexpr#1\columnwidth-2\tabcolsep-1.5\arrayrulewidth}}
\usepackage{booktabs}	% Kompatibilität für array und \hrule
\usepackage{caption}	% Um auch unbeschriftete Captions mit \caption* zu machen

% Nützliches
\usepackage{verbatim}	% u.a. zum auskommentieren via \begin{comment} \end{comment}
\usepackage{tabto}	% Tabs: /tab zum nächsten Tab oder /tabto{.5 \CurrentLineWidth} zur Stelle in der Linie
\NumTabs{6}	% Anzahl von Tabs pro Zeile zum springen
\usepackage{listings} % Source-Code mit Tabs
\usepackage{lstautogobble} 
\usepackage{enumitem}	% Anpassung der enumerates
%\setlist[enumerate,1]{label=(\arabic*)}	% global andere Enum-Items
\renewcommand{\labelitemiii}{$\scriptscriptstyle ^\blacklozenge$} % global andere 3. Item-Aufzählungszeichen
\usepackage{letltxmacro} % neue Definiton von Grundbefehlen
% Nutzen:
%\LetLtxMacro{\oldemph}{\emph}
%\renewcommand{\emph}[1]{\oldemph{#1}}
\RequirePackage{xpatch}	% ua. Konkatenieren von Strings/Variablen (etoolbox)
\usepackage{xstring}	% String Operationen
\usepackage{minibox}	% Minibox anstatt \fbox{} für Boxen mit (manuellen) Zeilenumbruch
\usepackage{framed}		% \begin{framed} statt minibox/fbox mit automatischem Zeilenumbruch
\usepackage{microtype}	% Verfeinerte Hyphenation
\usepackage{import}		
\usepackage{standalone}	% Zusammen mit import: Einfügen von separaten Latex-Dateien (die schon header haben)

% Einrichtung von lst
%\renewcommand*{\ttdefault}{qcr}
\lstset{
basicstyle=\ttfamily, 
%basicstyle=\footnotesize\sffamily\color{black},
%mathescape=true, 
%escapeinside=^^, 
autogobble, 
tabsize=2,
frame=single,
rulecolor=\color{lightgray},
numbers=left,
numbersep=5pt,
numberstyle=\tiny\color{gray},
commentstyle=\color{gray},
keywordstyle=\color{htwblue},
stringstyle=\color{htworange},
morecomment=[l][\color{magenta}]{\#}
showspaces=false,
showstringspaces=false,
breaklines=true,
literate=%
    {Ö}{{\"O}}1
    {Ä}{{\"A}}1
    {Ü}{{\"U}}1
    {ß}{{\ss}}1
    {ü}{{\"u}}1
    {ä}{{\"a}}1
    {ö}{{\"o}}1
    {„}{{\glqq}}1
    {“}{{\grqq}}1
    {~}{{\textasciitilde}}1
}
\usepackage{scrhack} % Fehler umgehen
\def\ContinueLineNumber{\lstset{firstnumber=last}} % vor lstlisting. Zum wechsel zum nicht-kontinuierlichen muss wieder \StartLineAt1 eingegeben werden
\def\StartLineAt#1{\lstset{firstnumber=#1}} % vor lstlisting \StartLineAt30 eingeben, um bei Zeile 30 zu starten
\let\numberLineAt\StartLineAt

% BibTeX
\usepackage[
backend=bibtex,
bibencoding=ascii,
style=authoryear, citestyle=authoryear-ibid,
%doi=false,
%isbn=false,
%url=false
]{biblatex}	% BibTeX
\usepackage{makeidx}
%\makeglossary
%\makeindex
\usepackage{url}	% Darstellung von URLs

% Grafiken
\usepackage{graphicx}
\usepackage{epstopdf}	% eps-Vektorgrafiken einfügen
\usepackage{transparent}	% transparent nutzen: {\transparent{0.4} ...}
%\epstopdfsetup{outdir=\customDir}
% Prüft, ob Grafik existiert (mit \ifvalidimage{}{}) [Quelle: https://tex.stackexchange.com/a/99176]:
\makeatletter
\newif\ifgraphicexist
\catcode`\*=11
\newcommand\ifvalidimage[1]{%
    \begingroup
    \global\graphicexisttrue
    \let\input@path\Ginput@path
    \filename@parse{#1}%
    \ifx\filename@ext\relax
    \@for\Gin@temp:=\Gin@extensions\do{%
        \ifx\Gin@ext\relax
        \Gin@getbase\Gin@temp
        \fi}%
    \else
    \Gin@getbase{\Gin@sepdefault\filename@ext}%
    \ifx\Gin@ext\relax
    \global\graphicexistfalse
    \def\Gin@base{\filename@area\filename@base}%
    \edef\Gin@ext{\Gin@sepdefault\filename@ext}%
    \fi
    \fi
    \ifx\Gin@ext\relax
    \global\graphicexistfalse
    \else 
    \@ifundefined{Gin@rule@\Gin@ext}%
    {\global\graphicexistfalse}%
    {}%
    \fi  
    \ifx\Gin@ext\relax 
    \gdef\imageextension{unknown}%
    \else
    \xdef\imageextension{\Gin@ext}%
    \fi 
    \endgroup 
    \ifgraphicexist
    \expandafter \@firstoftwo
    \else
    \expandafter \@secondoftwo
    \fi 
} 
\catcode`\*=12
\makeatother
\LetLtxMacro{\forceincludegraphics}{\includegraphics}	% neuer Befehl für includegraphics
\renewcommand{\includegraphics}[2][]{	% altes includegraphics neu definieren, damit es auch nicht vorhandene einfügt
\ifvalidimage{#2}{
\forceincludegraphics[#1]{#2}
}{
\message{Achtung: Grafik wurde nicht gefunden: '#2'}
\minibox[frame]{
\textbf{\StrSubstitute{#2}{_}{\_}}%  \ifthenelse{\isempty{#1}}{}{\\\textit{#1}}
}}}

% pdf-Setup
\usepackage{pdfpages}
\usepackage[bookmarks,%
bookmarksopen=false,% Klappt die Bookmarks in Acrobat aus
colorlinks=true,%
linkcolor=black,%
citecolor=black,%
urlcolor=htworange,%
]{hyperref}

% Nach Dokumentenbeginn ausführen:
\makeatletter
\AtBeginDocument{	
	% pdf Informationen
	\hypersetup{
		pdfauthor={\@author},
		pdftitle={\@title},
		pdfproducer={LaTeX Layout erstellt von Falk-Jonatan Strube}
	}
	% signierte Fußzeile
	%\providecommand{\ofootname}{Mitschrift von}
	\providecommand{\pathtomaster}{../}
	\ofoot{%
	%\ifdefined\ofoottext
	%\else
	%\begin{minipage}{40mm}
	%\begin{flushright}
	%\textcolor{gray}{
	%\ofootname\\
	%Falk-Jonatan Strube\vspace*{-1mm}}
	%\end{flushright}
	%\end{minipage}
	%\fi%
	%\hspace*{1mm}%
	\begin{minipage}{10 mm}
	\includegraphics[height=9mm]{\pathtomaster/FJS-logo_gray}
	\end{minipage}%
	}
}
\makeatother



%%% ==================================================
%%% ===================== MACROS =====================
%%% ==================================================
% Horizontale Linie:
\newcommand{\HRule}[1][\medskipamount]{\par
  \vspace*{\dimexpr-\parskip-\baselineskip+#1}
  \noindent\rule[0.2ex]{\linewidth}{0.2mm}\par
  \vspace*{\dimexpr-\parskip-.5\baselineskip+#1}}
% Gestrichelte horizontale Linie:
\RequirePackage{dashrule}
\newcommand{\HDRule}[1][\medskipamount]{\par
  \vspace*{\dimexpr-\parskip-\baselineskip+#1}
  \noindent\hdashrule[0.2ex]{\linewidth}{0.2mm}{1mm} \par
  \vspace*{\dimexpr-\parskip-.5\baselineskip+#1}}
% Mathe in Anführungszeichen:
\newsavebox{\mathbox}\newsavebox{\mathquote}
\makeatletter
\newcommand{\mq}[1]{% \mathquotes{<stuff>}
  \savebox{\mathquote}{\text{"}}% Save quotes
  \savebox{\mathbox}{$\displaystyle #1$}% Save <stuff>
  \raisebox{\dimexpr\ht\mathbox-\ht\mathquote\relax}{"}#1\raisebox{\dimexpr\ht\mathbox-\ht\mathquote\relax}{''}
}
\makeatother

% Paragraph mit Zähler (Section-Weise)
\newcounter{cparagraphC}
\newcommand{\cparagraph}[1]{
\stepcounter{cparagraphC}
\paragraph{\thesection{}-\thecparagraphC{} #1}
%\addcontentsline{toc}{subsubsection}{\thesection{}-\thecparagraphC{} #1}
\label{\thesection-\thecparagraphC}
}
\makeatletter
\@addtoreset{cparagraphC}{section}
\makeatother


% (Vorlesungs-)Folien einbinden:
\makeatletter
% Folien-Pfad (inkl. Dateiname ohne Endung und ggf. ohne Nummerierung)
\newcommand{\@slidePath}{}
\newcommand{\slidePath}[1]{\renewcommand{\@slidePath}{#1}}
% Folien Eigenschaften
\newcommand{\@slideScale}{0.5}
\newcommand{\slideScale}[1]{\renewcommand{\@slideScale}{#1}}
\AtBeginDocument{
% Folien von einer Datei skaliert
\newcommand{\slide}[2][\@slideScale]{\slides[#1]{}{#2}}
\newcommand{\slideTrim}[6][\@slideScale]{\slides[#1 , clip,  trim = #5cm #4cm #6cm #3cm]{}{#2}}
% Folien von mehreren nummerierten Dateien skaliert
\newcommand{\slides}[3][\@slideScale]{\begin{center}
\includegraphics[page=#3, scale=#1]{\@slidePath #2.pdf}
\end{center}}
}
\makeatother

% \emph{} anders definieren
\makeatletter
\DeclareRobustCommand{\em}{%
  \@nomath\em \if b\expandafter\@car\f@series\@nil
  \normalfont \else \scshape \fi}
\makeatother

% unwichtiges
\newcommand{\unimptnt}[1]{{\transparent{0.5}#1}}

% alph. enumerate
\newenvironment{anumerate}{\begin{enumerate}[label=(\alph*)]}{\end{enumerate}} % Alphabetische Aufzählung

% Hanging parameters
\newcommand{\hangpara}[1]{\par\noindent\hangindent+2em\hangafter=1 #1\par\noindent}

%% EINFACHE BEFEHLE

% Abkürzungen Mathe
\newcommand{\EE}{\mathbb{E}}
\newcommand{\QQ}{\mathbb{Q}}
\newcommand{\RR}{\mathbb{R}}
\newcommand{\CC}{\mathbb{C}}
\newcommand{\NN}{\mathbb{N}}
\newcommand{\ZZ}{\mathbb{Z}}
\newcommand{\PP}{\mathbb{P}}
\renewcommand{\SS}{\mathbb{S}}
\newcommand{\cA}{\mathcal{A}}
\newcommand{\cB}{\mathcal{B}}
\newcommand{\cC}{\mathcal{C}}
\newcommand{\cD}{\mathcal{D}}
\newcommand{\cE}{\mathcal{E}}
\newcommand{\cF}{\mathcal{F}}
\newcommand{\cG}{\mathcal{G}}
\newcommand{\cH}{\mathcal{H}}
\newcommand{\cI}{\mathcal{I}}
\newcommand{\cJ}{\mathcal{J}}
\newcommand{\cM}{\mathcal{M}}
\newcommand{\cN}{\mathcal{N}}
\newcommand{\cP}{\mathcal{P}}
\newcommand{\cR}{\mathcal{R}}
\newcommand{\cS}{\mathcal{S}}
\newcommand{\cZ}{\mathcal{Z}}
\newcommand{\cL}{\mathcal{L}}
\newcommand{\cT}{\mathcal{T}}
\newcommand{\cU}{\mathcal{U}}
\newcommand{\cX}{\mathcal{X}}
\newcommand{\cV}{\mathcal{V}}
\renewcommand{\phi}{\varphi}
\renewcommand{\epsilon}{\varepsilon}
\renewcommand{\theta}{\vartheta}

% Verschiedene als Mathe-Operatoren
\DeclareMathOperator{\arccot}{arccot}
\DeclareMathOperator{\arccosh}{arccosh}
\DeclareMathOperator{\arcsinh}{arcsinh}
\DeclareMathOperator{\arctanh}{arctanh}
\DeclareMathOperator{\arccoth}{arccoth} 
\DeclareMathOperator{\var}{Var} % Varianz 
\DeclareMathOperator{\cov}{Cov} % Co-Varianz 

% Farbdefinitionen
\definecolor{red}{RGB}{180,0,0}
\definecolor{green}{RGB}{75,160,0}
\definecolor{blue}{RGB}{0,75,200}
\definecolor{orange}{RGB}{255,128,0}
\definecolor{yellow}{RGB}{255,245,0}
\definecolor{purple}{RGB}{75,0,160}
\definecolor{cyan}{RGB}{0,160,160}
\definecolor{brown}{RGB}{120,60,10}

\definecolor{itteny}{RGB}{244,229,0}
\definecolor{ittenyo}{RGB}{253,198,11}
\definecolor{itteno}{RGB}{241,142,28}
\definecolor{ittenor}{RGB}{234,98,31}
\definecolor{ittenr}{RGB}{227,35,34}
\definecolor{ittenrp}{RGB}{196,3,125}
\definecolor{ittenp}{RGB}{109,57,139}
\definecolor{ittenpb}{RGB}{68,78,153}
\definecolor{ittenb}{RGB}{42,113,176}
\definecolor{ittenbg}{RGB}{6,150,187}
\definecolor{itteng}{RGB}{0,142,91}
\definecolor{ittengy}{RGB}{140,187,38}

% Textfarbe ändern
\newcommand{\tred}[1]{\textcolor{red}{#1}}
\newcommand{\tgreen}[1]{\textcolor{green}{#1}}
\newcommand{\tblue}[1]{\textcolor{blue}{#1}}
\newcommand{\torange}[1]{\textcolor{orange}{#1}}
\newcommand{\tyellow}[1]{\textcolor{yellow}{#1}}
\newcommand{\tpurple}[1]{\textcolor{purple}{#1}}
\newcommand{\tcyan}[1]{\textcolor{cyan}{#1}}
\newcommand{\tbrown}[1]{\textcolor{brown}{#1}}

% Umstellen der Tabellen Definition
\newcommand{\mpb}[1][.3]{\begin{minipage}{#1\textwidth}\vspace*{3pt}}
\newcommand{\mpe}{\vspace*{3pt}\end{minipage}}

\newcommand{\resultul}[1]{\underline{\underline{#1}}}
\newcommand{\parskp}{$ $\\}	% new line after paragraph
\newcommand{\parnl}{\unskip\parfillskip 0pt \par} % justified new line (brace line!)

% mathematische Abkürzungen
\newcommand{\corr}{\;\widehat{=}\;}
\newcommand{\mdeg}{^{\circ}}
\newcommand{\nok}[2]{\binom{#1}{#2}}	% n über k BESSER: \binom{n}{k}
\newcommand{\mtr}[1]{\begin{pmatrix}#1\end{pmatrix}}	% Matrix
\newcommand{\dtr}[1]{\begin{vmatrix}#1\end{vmatrix}}	% Determinante (Betragsmatrix)
\LetLtxMacro{\originalVec}{\vec}
\renewcommand{\vec}[1]{\underline{#1}}	% Vektorschreibweise
\newcommand{\imptnt}[1]{\colorbox{red!30}{#1}}	% Wichtiges
\newcommand{\intd}[1]{\,\mathrm{d}#1}
\newcommand{\diffd}[1]{\mathrm{d}#1}
% für Module-Rechnung: \pmod{}
\newcommand{\unit}[1]{\,\mathrm{#1}}
%\LetLtxMacro{\ntilde}{\tilde}
%\renewcommand{\tilde}{\widetilde}

% weiter Abkürzungen
\newcommand{\gdw}{genau dann wenn}
\newcommand{\lecdate}[1]{\marginpar{\small\textcolor{gray}{Vorlesung\newline #1}}}

\faculty{Fakultät Informatik/Mathematik}
\chair{Lehrstuhl für Programmierung, Compilerbau/Interpreter, Mikrocontroller}
\subject{lecture}
\title{Compiler/Interpreter}
\author{Falk-Jonatan Strube}
\professor{Prof. Dr.-Ing. Arnold Beck}
\slidePath{Vorlesung/}
\slideScale{.38}
\begin{document}
\maketitle
\tableofcontents

\chapter*{Hinweise}

Zugelassene Hilfsmittel Klausur: Spickzettel A-4 Blatt, doppelseitig\\
Wenn PL/0 compiler zum Ende des Semsters fertig: 10\% der Klausurpunkte Bonus.

\subsection*{Lehrinhalte}
\begin{itemize}
\item Einführung in die Begriffswelt der Theorie der formalen Sprachen
\item Strategien der Analyse
\item Aufbau von Compilern/Interpretern
\item Arbeitsweise höherer Programmiersprachen
\item Arbeiten mit Automaten und Graphen und deren programmiertechnische Umsetzung
\end{itemize}

\chapter{Einführung}
\section{Motivation}
\slides{V1_Grundlagen}{5}
\subsection{Der Compiler}
\slides{V1_Grundlagen}{6}
\subsubsection{Aufbau eines Compilers}
\slides{V1_Grundlagen}{7}
\subsubsection{Beispiel}
\slides{V1_Grundlagen}{8}
\subsubsection*{Abstrakter Syntaxbaum}
\slides{V1_Grundlagen}{9}
\slides{V1_Grundlagen}{10}
\subsection{Dreiadressmaschine vs Stackmaschine}
\slides{V1_Grundlagen}{11}
\slides{V1_Grundlagen}{12}
\subsection{Syntaxgesteuerter Einpasscompiler}
\slides{V1_Grundlagen}{13}

\section{Grundlagen}
\slides{V2_Ableitung}{1}
\subsection{Sprache}
\slides{V1_Grundlagen}{14}
\subsection{Alphabet}
\slides{V1_Grundlagen}{15}
\subsubsection*{Beispiel}
\slides{V1_Grundlagen}{16}
\subsection{Zeichenkette}
\slides{V1_Grundlagen}{17}
\subsection{Freie Halbgruppe (A*)}
\slides{V1_Grundlagen}{18}
\subsection{Syntaxregeln}
\slides{V1_Grundlagen}{18}
\subsection{Grammatik und Parser}
\slides{V2_Ableitung}{6}
\subsection{BNF}
\slides{V2_Ableitung}{7}
\slides{V2_Ableitung}{8}
\subsubsection*{Beispiel}
\slides{V2_Ableitung}{9}
\subsection{EBNF}
\slides{V2_Ableitung}{10}
\slides{V2_Ableitung}{11}
\subsection{Syntaxgraph}
\slides{V2_Ableitung}{12}
\slides{V2_Ableitung}{13}
\slides{V2_Ableitung}{14}
\slides{V2_Ableitung}{15}
\subsection{Grammatik von PL/0}
\slides{V2_Ableitung}{16}
\slides{V2_Ableitung}{17}
\subsubsection*{Beispiele}
\slides{V2_Ableitung}{18}

\section{Grammatik nach Chomsky}
\slides{V1_Grundlagen}{20}
\subsubsection*{Beispiel}
\slides{V1_Grundlagen}{21}
\subsection{Ableitung}
\slides{V1_Grundlagen}{22}
\subsubsection*{Direkte Ableitung}
\slides{V1_Grundlagen}{23}
\slides{V1_Grundlagen}{24}
\subsection{Chomsky Typ 0}
\slides{V1_Grundlagen}{25}
\subsection{Chomsky Typ 1}
\slides{V1_Grundlagen}{26}
\subsection{Chomsky Typ 2}
\slides{V1_Grundlagen}{27}
\subsection{Chomsky Typ 3}
\slides{V1_Grundlagen}{28}
\subsection{Chomsky Typ 4}
\slides{V1_Grundlagen}{29}
\subsection{Beispiel}
\slides{V1_Grundlagen}{30}
\subsubsection*{Analysen}
\slides{V1_Grundlagen}{31}
\slides{V1_Grundlagen}{32}
\subsection{Analysestrategien}
\slides{V1_Grundlagen}{33}

\section{Morpheme, Token, Atome}
\slides{V1_Grundlagen}{34}
\subsection{Morpheme einfacher Ausdrücke}
\slides{V1_Grundlagen}{35}
\slides{V1_Grundlagen}{36}
\subsubsection*{Beispiel 3+5*12}
\slides{V1_Grundlagen}{37}
\subsection{Ein einfacher Lexer}
\begin{lstlisting}[language=C]
typedef struct morph{
	int mc;
	double dval;
	char cval;
}MORPHEM;

enum mcodes{
	mempty,
	mop,
	mdbl
};

static MORPHEM m;

void lex(char * pX){
	static char * px;
	/* Initialisierung*/
	if (pX!=NULL){
		m.mc=mempty;
		px=pX;
	}
	/* lexiaklische Analyse */
	if (*px=='\0'){
		m.mc=mempty;
	} else {
		for (;*px==' '||*px=='\t';px++);
		if (isdigit(*px) || *px=='.'){
			m.dval=strtod(px,&px);
			m.mc =mdbl;
		} else
			switch (*px){
				case '+':
				case '':
				case '*':
				case '/':
				case '(':
				case ')':
					m.cval=*px++;
					m.mc=mop;
					break;
				default :
					printf("wrong ...: %c\n-- canceling --\n",*px);
					exit (1);
			}
	}
}
\end{lstlisting}

\section{Grammatik 1: Einfache Ausdrücke}
\slides{V1_Grundlagen}{39}
\subsection{Rekursiver Abstieg}
\slides{V1_Grundlagen}{40}
\subsubsection*{Alternative Regeln}
\begin{lstlisting}[language=C]
if (lookahead == First(A1)){
	A1();
} else if (lookahead == First(A2)) {
	A2();
}
/* ... */
else if (lookahead == First(An)){
	An();
} else {
	error();
}
\end{lstlisting}
First ist das jeweils erste terminale Symbol einer alternativen
Regel
\subsection{Anforderungen}
\slides{V1_Grundlagen}{42}

\section{LL(1) Grammatiken}
\slides{V1_Grundlagen}{43}
\slides{V1_Grundlagen}{44}

\section{Grammatik 2: Einfache Ausdrücke ohne Linksrekursion}
\slides{V1_Grundlagen}{45}
\subsection{Linksfaktorisierung}
\slides{V1_Grundlagen}{46}
\subsubsection*{An Grammatik}
\slides{V1_Grundlagen}{47}
\subsection{First/Follow}
\slides{V1_Grundlagen}{48}
\subsection{Expression}
\begin{lstlisting}[language=C]
double expr(void){
	double tmp=term();
	if (m.mc==mop && m.cval=='+'){
		lex(NULL); 
		tmp+=expr();
	}
	return tmp;
}
\end{lstlisting}
\subsection{Term}
\begin{lstlisting}[language=C]
double term(void){
	double tmp=fac();
	if (m.mc==mop && m.cval=='*'){
		lex(NULL);
		tmp*=term();
	}
	return tmp;
}
\end{lstlisting}
\subsection{Faktor}
\begin{lstlisting}[language=C]
double fac(){
	double tmp;
	if (m.mc==mop)	{
		if (m.cval=='(')	{
			lex(NULL);
			tmp=expr();
			if (m.mc != mop || m.cval != ')')
				exit (1);
			lex(NULL);
		} else 
			exit (1);
	}	else
		if (m.mc==mdbl){
			tmp=m.dval;
			lex(NULL);
		}	else 
			exit (1);
	return tmp;
}
\end{lstlisting}
\subsection{Main}
\begin{lstlisting}[language=C]
int main(int argc, char*argv[]){
	char *pBuf=argv[1];
	printf("%s\n",pBuf);
	lex(pBuf);
	printf("%10.4f\n",expr());
	free(pBuf);
	return 0;
}
\end{lstlisting}
\subsection{Ausführen}
\begin{lstlisting}
beck@linux:~/COMPILER> ./tPM 5+2*3
5+2*3
11.0000
aber:
beck@linux:~/COMPILER>./tPM 12/2*3
\end{lstlisting}
$\to$ Grammatik rechnet nicht korrekt!

\section{Umformung der Grammatik}
\lecdate{25.10.2017}
\slides{V1_Grundlagen}{54}

\subsection{Grammatik 3: Beseitigung der Linksrekursion durch Iteration}
\slides{V1_Grundlagen}{55}

\subsection{EBNF}
\slides{V1_Grundlagen}{56}

\subsection{Expression als C-Code}
\slides{V1_Grundlagen}{57}

\subsection{Beseitigung der Linksrekursion bei klassischer BNF}
\slides{V1_Grundlagen}{58}

\section{Grammatik für einfache Ausdrücke}
\slides{V1_Grundlagen}{59}

\section{First/Follow}
\slides{V1_Grundlagen}{60}
\subsection{Implementation}
\slides{V1_Grundlagen}{61}
\subsubsection*{Syntaxgraph}
\slides{V1_Grundlagen}{62}

\section{Wiederholung}
\slides{V1_Grundlagen}{63}








%\chapter{Ableitung} %Integriert in Einführung
%\slides{V2_Ableitung}{2}

\chapter{Lexer}
\section{Lexikalische Analyse}
\slides{V3_Lexer}{1}

\subsection{Einordnung in den Compiler}
\slides{V3_Lexer}{2}

\subsection{Funktionsweise}
\slides{V3_Lexer}{3}

\section{Grammatik der PL/0 Token}
\slides{V3_Lexer}{4}
Buzi: Buchstabe/Ziffer; Bu: Buchstabe
\subsection{Regeln der Grammatik der Morpheme}
\slides{V3_Lexer}{5}
\subsubsection*{Syntaxgraph}
\label{pl0syntaxgraph}
\slides[.4]{V3_Lexer}{6}

\section{Endlicher deterministischer Automat}
\slides{V3_Lexer}{7}
\slides{V3_Lexer}{8}
\slides{V3_Lexer}{9}

\subsection{Automatentabelle}
\slides{V3_Lexer}{10}
\subsubsection{Zeichenklassenvektor}
\slides{V3_Lexer}{11}
\slides[.4]{V3_Lexer}{12}
$\to$Jedem ASCII-Zeichen wird eine Spalte („Handlung“) in der Automatentabelle zugewiesen (mit Hilfe des Syntaxgraphs in \autoref{pl0syntaxgraph}):
\slides[.4]{V3_Lexer}{15}
s: Schreibe des betrachteten Symbols in die Ausgabe\\
g/sg: Schreibe Großbuchstabe des betrachteten Buchstaben\\
l: Lese nächstes Symbol\\
b: Beende das Erkennen dieses Tokens

\subsubsection{Aktionen}
\slides{V3_Lexer}{16}

\subsubsection{Fazit}
\slides{V3_Lexer}{17}

\subsubsection*{Variablen Beispiellexer}
\slides{V3_Lexer}{18}

\section{Funktionen des Lexers}
\slides{V3_Lexer}{19}
\subsection{Beispielimplementierung}
\slides{V3_Lexer}{20}
\slides[.4]{V3_Lexer}{21}
%\slides{V3_Lexer}{22}
%\slides{V3_Lexer}{23}
\begin{lstlisting}[language=C]
/* Schalt und Ausgabefunktionen des Automaten */
static void fl (void);
static void fb (void);
static void fgl (void);
static void fsl (void);
static void fslb(void);
typedef void (*FX)(void);
static FX vfx[]={fl,fb,fgl,fsl,fslb};

/* Funktionsindex *0x10, bzw. *16 */
typedef enum T_Fx{
	ifl=0x0,
	ifb=0x10,
	ifgl=0x20,
	ifsl=0x30,
	ifslb=0x40
}tFx;

/* Morphemcodes */
typedef enum T_MC{
	mcEmpty, mcSymb, mcNum, mcIdent
}tMC;
typedef enum T_ZS{
	zNIL,
	zERG=128,zLE,zGE,
	zBGN,zCLL,zCST,zDO,zEND,zIF,zODD,zPRC,zTHN,zVAR,zWHL
}tZS;

typedef struct{
	tMC MC; /* Morphemcode */
	int PosLine; /* Zeile */
	int PosCol; /* Spalte */
	union VAL	{
		long Num;
		char*pStr;
		int Symb;
	}Val;
	int MLen; /* Morphemlnge*/
}tMorph;

int initLex(char* fname);
tMorph* Lex(void);
\end{lstlisting}
\slides[.4]{V3_Lexer}{24}
\slides[.4]{V3_Lexer}{25}
\begin{lstlisting}[language=C]
static FILE *pIF;
static tMorph MorphInit;
extern tMorph Morph;
static int X;
static int Z;
static char vBuf[1024+1];
static char* pBuf;
static int line,col;
static int Ende; // Entfällt, wenn Zustand 9 -> Ende
/* Initialisierung der lexiaklischen Analyse */
int initLex(char* fname){
	char vName[128+1];
	strcpy(vName,fname);
	if (strstr(vName,".pl0")==NULL)
		strcat(vName,".pl0");
	pIF=fopen(vName,"r+t");
	if (pIF!=NULL){
		X=fgetc(pIF); /* Lesen des ersten Zeichens */
		return OK;
	}
	return FAIL;
}
\end{lstlisting}

\subsection{Algorithmus des Lexers}
\begin{lstlisting}[language=C]
tMorph* Lex(void){
	Z=0; // Anfangszustand, wenn nicht aus Tabelle
	int zx;
	Morph=MorphInit;
	Morph.PosLine=line;
	Morph.PosCol =col;
	pBuf=vBuf;
	End=0;
	do{
		zx= vSMatrix[Z][vZKl[X&0x7f]]&0xF;
		vfx[(vSMatrix[Z][vZKl[X&0x7f]])>>4]();
		Z=zx;
	}while (End==0); // (Z!=zEnd) // zEnd:9
	return &Morph;
}
\end{lstlisting}

Hinweis Praktikum: Klassifizierung soll über Zeichenklassenvektor passieren!










\chapter{Parser}
\section{Grundlagen}
\slides{V4_Parser}{1}
\subsection{Kellerautomat}
\slides{V4_Parser}{2}
\slides{V4_Parser}{3}

\section{LL(1)-Grammatik}
\slides{V4_Parser}{4}
\slides{V4_Parser}{5}
Hinweis: In PL0 darf das \lstinline`statement` auch leer sein (da eckige Klammern).
\subsection{First/Follow}
\slides{V4_Parser}{6}
\subsection{Als Graph}
\subsubsection{programm}
\slides{V4_Parser}{7}
\subsubsection{block}
\slides{V4_Parser}{8}
\subsubsection{statement}
\slides{V4_Parser}{9}
\subsubsection{expr}
\lecdate{08.11.2017}
\slides{V4_Parser}{10}
1. Variante reicht für Parser, 3. Variante ist sinnvoller, wenn das Ergebnis im Compiler weiter genutzt werden soll.
\subsubsection{term}
\slides{V4_Parser}{11}
\subsubsection{factor}
\slides{V4_Parser}{12}
\subsubsection{condition}
\slides{V4_Parser}{13}
\subsubsection{Bewertung}
\slides{V4_Parser}{14}
\slides{V4_Parser}{15}
\subsubsection*{Alternative}
\slides{V4_Parser}{16}
\slides{V4_Parser}{17}

\section{Implementation von Graphen}
\slides{V4_Parser}{18}
\subsection{Struktur eines Bogens}
\slides{V4_Parser}{19}
\subsubsection*{Beispiel Factor}
\begin{lstlisting}[language=C]
typedef unsigned long ul;
tBog gFact[] = {
/* 0*/ {BgMo,{(ul)mcIdent	}, NULL, 5, 1},	/*(0)---ident--->(E)*/
/* 1*/ {BgMo,{(ul)mcNumb 	}, NULL, 5, 2},	/* +----number-->(E)*/
/* 2*/ {BgSy,{(ul)'(' 		}, NULL, 3, 0},	/*(+)----'('---->(3)*/
/* 3*/ {BgGr,{(ul)gExpr 	}, NULL, 4, 0},	/*(1)--express-->(4)*/
/* 4*/ {BgSy,{(ul)')' 		}, NULL, 5, 0},	/*(0)----')'---->(E)*/
/* 5*/ {BgEn,{(ul)0 			}, NULL, 0, 0} 	/*(E)---------(ENDE)*/
};
\end{lstlisting}
\slides[.5]{V4_Parser}{20}
Hinweis: $0$ in \lstinline`iAlt` bedeutet keine Alternative (vorher festgelegt, könnte auch $-1$ sein).
\subsection{Erstellen eines Graphen}
\slides[.5]{V4_Parser}{21}
\begin{lstlisting}[language=C]
tBog gFact[]={
/* 0*/ {BgMo,{(ul)mcIdent	}, NULL, 5, 1},
/* 1*/ {BgMo,{(ul)mcNumb	}, NULL, 5, 2},
/* 2*/ {BgSy,{(ul)'(' 		}, NULL, 3, 0},
/* 3*/ {BgGr,{(ul)gExpr 	}, NULL, 4, 0},
/* 4*/ {BgSy,{(ul)')' 		}, NULL, 5, 0},
/* 5*/ {BgEn,{(ul)0 			}, NULL, 0, 0}
};
\end{lstlisting}
\subsubsection*{Hinweise zur Umsetzung}
\slides{V4_Parser}{22}
\subsection{Implementation von Bögen in Graphen}
\slides{V4_Parser}{23}
\begin{lstlisting}[language=C]
int pars(tBog* pGraph){
	tBog* pBog=pGraph;
	int succ=0;
	if (Morph.MC==mcEmpty)Lex();
	while(1){
		switch(pBog->BgD){
			case BgNl:
				succ=1;
				break;
			case BgSy:
				succ=(Morph.Val.Symb==pBog->BgX.S);
				break;
			case BgMo:
				succ=(Morph.MC==(tMC)pBog->BgX.M);
				break;
			case BgGr:
				succ=pars(pBog->BgX.G);
				break;
			case BgEn:
				return 1; /* Ende erreichet Erfolg */
		}
		if (succ && pBog->fx!=NULL)
			succ=pBog->fx();
		if (!succ)/* Alternativbogen probieren */
			if (pBog->iAlt != 0)
				pBog=pGraph+pBog->iAlt;
			else
				return FAIL;
		else{/* Morphem formal akzeptiert (eaten) */
			if (pBog->BgD & BgSy || pBog->BgD & BgMo)
				Lex();
			pBog=pGraph+pBog->iNext;
		}
	}/* while */
}
\end{lstlisting}
\subsection{Parsebaum (optional)}
\slides{V4_Parser}{25}
\begin{lstlisting}[language=C]
typedef struct{
	int 		line;
	char 		Type[10];
	tList* 	pList;
	char * 	Descr;
}TreeItem;

TreeItem* crItem(	tList* pList, // neue Liste od. NULL
									char* Descr, // Description
									char* Type, // Bogentyp
									int line){ // Quellzeile
	TreeItem *ptmp=malloc(sizeof(TreeItem));
	strncpy(ptmp->Type,Type,9);
	ptmp->pList=pList;
	ptmp->Descr=Descr;
	ptmp->line=line;
	return ptmp;
}
\end{lstlisting}
Aufruf:
\begin{lstlisting}[language=C]
InsertHead(pTree,crItem((pl=CreateList()),"Prog","MetaSymb",0));
if (pars(gProg,pl)==1)
/* ... */

int pars(tBog* pGraph,tList* pList){
	int verarbMorph=0;
	tBog* pBog=pGraph;
	int succ;
	tList* pl=NULL;
	TreeItem *pItem=NULL;
	if (Morph.MC==mcEmpty)Lex();
	while(1){
		switch(pBog->BgD & (BgNl+BgSy+BgMo+BgGr+BgEn)){
			case BgNl:
				succ=1;
				break;
			case BgSy:
				succ=(Morph.Val.Symb==pBog->BgX.S);
				if (succ)
					InsertTail(pList,crItem(NULL,(crStr(Morph.vBuf)),"Symbol",Morph.PosLine));
				break;
			case BgMo:
				succ=(Morph.MC==(tMC)pBog->BgX.M);
				if (succ)
					InsertTail(pList,crItem(NULL,(crStr(Morph.vBuf)),"Morph",Morph.PosLine));
				break;
			case BgGr:
				pl=CreateList();
				pItem=crItem((pl),crStr(StrGr[pBog->BgX.G]),"MetaSymb",Morph.PosLine);
				InsertTail(pList,pItem);
				succ=pars(vGr[pBog->BgX.G],pl);
				if (succ!=OK){
					DeleteList(pl); 
					free(pItem->Descr);
					free(pItem); 
					GetLast(pList); 
					RemoveItem(pList);
				}
				break;
			case BgEn:
				return 1; /* Ende erreichet - Erfolg */
		}
		
		if (succ && pBog->fx!=NULL)
			succ=pBog->fx();
		if (!succ){/* Alternativbogen probieren */
			if (pBog->iAlt != 0)
				pBog=pGraph+pBog->iAlt;
			/* Graph nicht erfolgreich verlassen */
			else if (verarbMorph)
				return ERROR; /* Syntaxfehler */
			else{
				//printf("go back\n");
				while(pItem=GetFirst(pList)){
					free(pItem->Descr);
					free(pItem);
					RemoveItem(pList);
				}
				return FAIL; /* vielleicht gibt es noch Alternative */
			}
		} else { /* Morphem formal akzeptiert */
			if (pBog->BgD & BgSy || pBog->BgD & BgMo){
				verarbMorph=1;
				Lex();
			}
			pBog=pGraph+pBog->iNext;
		}
	}
}
\end{lstlisting}
\slides{V4_Parser}{27}

\subsection{Erweiterte Graphen mit Backtracking}
\slides[.5]{V4_Parser}{30}
\slides[.5]{V4_Parser}{31}
\slides[.5]{V4_Parser}{32}
\slides[.5]{V4_Parser}{33}
\begin{lstlisting}[language=C]
tBog gStmnt[]= { 
/* 0*/ {BgGr,{(unsigned long)iAssign	}, NULL, 7, 1},
/* 1*/ {BgGr,{(unsigned long)iCall 		}, NULL, 7, 2},
/* 2*/ {BgGr,{(unsigned long)iBegin 	}, NULL, 7, 3},
/* 3*/ {BgGr,{(unsigned long)iIf 			}, NULL, 7, 4},
/* 4*/ {BgGr,{(unsigned long)iWhile 	}, St4 , 7, 5},
/* 5*/ {BgGr,{(unsigned long)iInput 	}, NULL, 7, 6},
/* 6*/ {BgGr,{(unsigned long)iOutput	}, NULL, 7, 0},
/* 7*/ {BgEn,{(unsigned long)0 				}, NULL, 0, 0}
};
tBog gAssign[]={
/* 0*/ {BgMo,{(unsigned long)mcIdent	}, St0 , 1, 0},
/* 1*/ {BgSy,{(unsigned long)zErg 		}, NULL, 2, 0},
/* 2*/ {BgGr,{(unsigned long)iExpr 		}, St8 , 3, 0},
/* 3*/ {BgEn,{(unsigned long)0 				}, NULL, 0, 0}
};
\end{lstlisting}







\chapter{Namensliste}
\section{Grundlagen}
\slides{V5_NamenList}{1}
\subsection*{Implementationsvarianten}
\slides{V5_NamenList}{2}
\section{Namenslisteneintrag}
\begin{lstlisting}[language=C]
typedef struct{
	tKz Kz;					// Kennzeichen (enum)
	short IdxProc;	// Prozedurnummer
	void* pObj;			// Poiner auf benanntes Objekt
	int Len;				
	char* pName;		// Pointer auf Name
}tBez;

typedef enum tKZ{
	KzBez, 
	KzPrc, 
	KzConst, 
	KzVar, 
	KzLabl, 
	Kz0ptr,
	// ...
}
\end{lstlisting}
\section{Variablenbeschreibung}
\slides{V5_NamenList}{4}
\begin{lstlisting}[language=C]
typedef struct tVAR{
	tKz Kz;
	int Dspl;	// Displacement
}tVar;
\end{lstlisting}
\section{Konstantenbeschreibung}
\slides{V5_NamenList}{6}
\begin{lstlisting}[language=C]
typedef struct tCONST{
	tKz Kz;
	long Val;	// Wert der Konstanten
	int Idx;	// Index der Konstanten im ConstBlock
}tConst;
\end{lstlisting}
%\subsubsection*{Konstanten am Ende des Codes}
%\slides{V5_NamenList}{7}
\section{Prozedurbeschreibung}
\slides{V5_NamenList}{9}
\begin{lstlisting}[language=C]
typedef struct tPROC{
	tKz Kz;								
	short IdxProc;				// Prozedurnummer
	struct tPROC*pParent;	// Pointer auf die Prozedurbeschreibung der Parentprozedur
	tList *pLBez;					// Lokale Namensliste der Prozedur
	int SpzzVar;					// Speicherplatzzuordnungszähler für Variable
}tProc;
\end{lstlisting}

\section{Beispiel Namensliste}
\slides[.5]{V5_NamenList}{11}
\begin{center}
$\downarrow$
\end{center}
\slides[.5]{V5_NamenList}{12}
\subsection*{Realisierung als Stack}
\slides[.5]{V5_NamenList}{13}

\section{Funktionen zur Namensliste}
\slides{V5_NamenList}{14}
\section{Einordnung in Parser}
bl* wird ausgeführt, wenn Bogen formal akzeptiert wurde.
\slides[.5]{V5_NamenList}{15}
\slides{V5_NamenList}{16}
\slides{V5_NamenList}{17}
\slides[.5]{V5_NamenList}{18}








\chapter{Virtuelle Maschine}
\section{Grundlagen}
\lecdate{29.11.2017}
\slides{V7_VirtuelleMasch}{2}

\subsection{Dreiadressmaschine}
\slides{V7_VirtuelleMasch}{3}
\subsection{Zweiadressmaschine}
\slides{V7_VirtuelleMasch}{4}
\subsection{Einadressmaschine}
\slides{V7_VirtuelleMasch}{5}
\subsection{Nulladressmaschine - Stackmaschine}
\slides{V7_VirtuelleMasch}{6}

\section{Befehlssatz}
\slides{V7_VirtuelleMasch}{7}
\begin{footnotesize}
\begin{lstlisting}[language=C]
typedef enum TCODE{
/*--- Kellerbefehle ---*/
puValVrLocl,/*00 (short Displ)	[Kellern Wert lokale Variable]								*/
puValVrMain,/*01 (short Displ)	[Kellern Wert Main Variable]									*/
puValVrGlob,/*02 (short Displ,short Proc) [Kellern Wert globale Variable]			*/
puAdrVrLocl,/*03 (short Displ)	[Kellern Adresse lokale Variable]							*/
puAdrVrMain,/*04 (short Displ)	[Kellern Adresse Main Variable] 							*/
puAdrVrGlob,/*05 (short Displ,short Proc) [Kellern Adresse globale Variable]	*/
puConst,		/*06 (short Index)	[Kellern einer Konstanten]										*/
storeVal,		/*07 ()							[Speichern Wert ->Adresse, beides aus Keller]	*/
putVal,			/*08 ()							[Ausgabe eines Wertes aus Keller nach stdout]	*/
getVal,			/*09 () 		[Eingabe eines Wertes von stdin -> Addr. im Keller ]	*/
/*--- arithmetische Befehle ---*/
vzMinus,		/*0A ()						 	[Vorzeichen ]																	*/
odd,				/*0B ()						 	[ungerade -> 0/1]																*/
/*--- binäre Operatoren kellern 2 Operanden aus und das Ergebnis ein ---*/
OpAdd,			/*0C ()						 	[Addition] 																		*/
OpSub,			/*0D ()						 	[Subtraktion ] 																*/
OpMult,			/*0E ()						 	[Multiplikation ] 														*/
OpDiv,			/*0F ()						 	[Division ] 																	*/
cmpEQ,			/*10 ()						 	[Vergleich =  -> 0/1] 												*/
cmpNE,			/*11 ()						 	[Vergleich #  -> 0/1] 												*/
cmpLT,			/*12 ()						 	[Vergleich <  -> 0/1] 												*/
cmpGT,			/*13 ()						 	[Vergleich >  -> 0/1] 												*/
cmpLE,			/*14 ()						 	[Vergleich <= -> 0/1]													*/
cmpGE,			/*15 ()						 	[Vergleich >= -> 0/1] 												*/
/*--- Sprungbefehle ---*/
call,				/*16 (short ProzNr) [Prozeduraufruf]															*/
retProc,		/*17 ()							[Rücksprung]																	*/
jmp,				/*18 (short RelAdr) [SPZZ innerhalb der Funktion]									*/
jnot,				/*19 (short RelAdr) [SPZZ innerhalb der Funkt.,Beding.aus Keller]	*/
entryProc,	/*1A (short lenCode,short ProcIdx,short lenVar)										*/
/*--- Zusätzliche, (für uns) nicht benötigte Befehle ---*/
putStrg,		/*1B (char[]) 																										*/
pop,				/*1C																															*/
swap,				/*1D								Austausch Adresse gegen Wert									*/
EndOfCode 	/*1E 																															*/
} tCode;
\end{lstlisting}
\end{footnotesize}
\subsection{Arbeitsweise der Stackmaschine}
\slides{V7_VirtuelleMasch}{8}

\section{Datenstrukturen der VM}
\slides{V7_VirtuelleMasch}{9}

\section{Register der VM}
\slides{V7_VirtuelleMasch}{10}

\section{Prozedurtabelle}
\slides{V7_VirtuelleMasch}{11}

\section{Initialisierung}
\slides{V7_VirtuelleMasch}{12}
\slides{V7_VirtuelleMasch}{13}

\section{Steuerschleife}
\begin{lstlisting}[language=C]
typedef int (*fx)(void);
fx vx[]={
	FcpuValVrLocl, FcpuValVrMain,
	FcpuValVrGlob, FcpuAdrVrLocl,
	FcpuAdrVrMain, FcpuAdrVrGlob,
	FcpuConst, FcstoreVal,
	FcputVal, FcgetVal,
	FcvzMinus, Fcodd,
	FcOpAdd, FcOpSub,
	FcOpMult, FcOpDiv,
	FccmpEQ, FccmpNE,
	FccmpLT, FccmpGT,
	FccmpLE, FccmpGE,
	Fccall, FcRetProc,
	Fcjmp, Fcjnot,
	FcEntryProc, FcputStrg,
	Fpop, Fswap,
	FcEndOfCode, Fcput,
	Fcget
};

while (!Ende){
	vx[*pC++]();
}

while (!Ende){
	switch (*pC++){
		case puValVrLocl:FcpuValVrLocl(); break;
		case puValVrMain:FcpuValVrMain(); break;
		case puValVrGlob:FcpuValVrGlob(); break;
		case puAdrVrLocl:FcpuAdrVrLocl(); break;
		case puAdrVrMain:FcpuAdrVrMain(); break;
		case puAdrVrGlob:FcpuAdrVrGlob(); break;
		case puConst :FcpuConst(); break;
		case storeVal :FcstoreVal(); break;
		case putVal :FcputVal(); break;
		case getVal :FcgetVal(); break;
		case vzMinus :FcvzMinus(); break;
		case odd :Fcodd(); break;
		case OpAdd :FcOpAdd(); break;
		case OpSub :FcOpSub(); break;
		case OpMult :FcOpMult(); break;
		case OpDiv :FcOpDiv(); break;
		/* ... */
		case pop :Fpop(); break;
		case swap :Fswap(); break;
		case EndOfCode :FcEndOfCode(); break;
		case put :Fcput(); break;
		case get :Fcget(); break;
	}
}
\end{lstlisting}

\section{Prozeduraufruf}
\slides{V7_VirtuelleMasch}{16}
\begin{lstlisting}[language=C]
int Fccall(void){
	short ProcNr;
	ProcNr=gtSrtPar(pC);pC+=2;
	push4((int4_t)(pInfProc[iCProc].pVar));
	push4((int4_t)pC);
	pC=pInfProc[ProcNr].pCode;//Sprung zu entryProc
	push4(iCProc);
	return OK;
}
int FcRetProc(void){
	pS=pInfProc[iCProc].pVar;
	iCProc =pop4();
	pC =(char*)pop4();
	pInfProc[iCProc].pVar=(char*)pop4();
	return OK;
}
\end{lstlisting}
\slides{V7_VirtuelleMasch}{18}
\begin{lstlisting}[language=C]
/* Entry Procedure */
int FcEntryProc(void){
	short lVar;
	/* Codelaenge ueberlesen */
	PC+=2;
	/* Prozedurnummer lesen und einstellen */
	iCProc=gtSrtPar(pC);pC+=2;
	/* Variablenlaenge lesen */
	lVar =gtSrtPar(pC);pC+=2;
	/* Zeiger auf Variablenbereich */
	pInfProc[iCProc].pVar=pS;
	/* Neuer Stackpointer */
	pS+=lVar;
	return OK;
}
\end{lstlisting}

\section{Rekursion}
\slides{V7_VirtuelleMasch}{20}

\section{Prozedur mit Parametern}
Parameter wird vor Aufruf auf den Stack geschrieben. Dieser hat die relative Adresse -16 (die erste lokale Variable hat die relative Adresse 0, Die Adresse der Funktion sind 3 Byte, daher haben dann Parameter die Adresse ab -16).
\slides{V7_VirtuelleMasch}{21}



\chapter{Codegenerierung}
%\section{Grundlagen}
%\slides{V6_Codegen}{1}













\chapter{Beispielcompilierung}
\lecdate{13.12.2017}

\section{Anzulegende Speicherbereiche}
\begin{itemize}
\item Morphem
\item ProcMain (Prozedurbeschreibung der Main Funktion)
\item Speicherbereich für Konstanten
\item Speicherbereich für Zwischencode
\item Ausgabedatei
\end{itemize}

\section{Abarbeitung}
\begin{itemize}
\item Lexer erkennt Morphem
\item Parser startet mit Startbogen
\item Parser ruft Blockbogen auf (rekursiver Aufruf von parse)
\item Vor dem erstellen einer Konstanten/Variablen: Prüfen, ob es schon einen mit gleichen Namen gibt.
\item Einfügen einer Konstante: erst prüfen, ob Wert bereits vorhanden. Konstanten mit gleichem Wert zeigen auf gleichen Speicherbereich
\item Zuweisung einer Variablen: pusht Konstante. D.h. Konstantenspeicherbereich wird ggf. erweitert, wenn Wert noch nicht drin (als nicht benannter Eintrag)
\end{itemize}







\end{document}